% TOP_PROBLEM: {"id":30007630,"problem_number":"LOCAL-30007630","title":"Commodity-trade fragmentation","category_id":21,"category":{"id":21,"name":"economics","display_name":"Economics","description":"Open economic theory statements, published research agendas, and proposed scoped specifications; question provenance and historical priority are qualified per record.","slug":"economics","order_index":21,"created_at":"2026-10-08T00:00:00Z"},"status":"research_question","external_url":null,"legacy_ids":[],"published":false,"statement":"In the explicitly specified setting: How do restrictions on food, energy and critical-mineral trade change scarcity, substitution and development risks? The mathematical statement above fixes the target, assumptions and scope of this catalogue formulation.","economics":{"original_id":119,"field":"Trade, globalization and geoeconomics","kind":"empirical","formalization_basis":"adapted_research_question","source_status":"research_agenda","priority_status":"earliest_found","priority_note":"Earliest verified explicit extension agenda located in this audit; no first-statement priority claim for commodity fragmentation.","earliest_verified_reference":{"authors":["Jorge A. Alvarez","Mehdi Benatiya Andaloussi","Chiara Maggi","Alexandre Sollaci","Martin Stuermer","Petia Topalova"],"title":"Geoeconomic Fragmentation and Commodity Markets","year":2023,"url":"https://www.imf.org/-/media/files/publications/wp/2023/english/wpiea2023201-print-pdf.pdf","doi":"10.5089/9798400252426.001","locator":"Section 10, concluding research-agenda paragraph, PDF pages 39–40","evidence_summary":"Requests critical-mineral, food-security, recycling and elasticity extensions beyond simplifying scenarios."},"current_reference":{"authors":["Jorge A. Alvarez","Mehdi Benatiya Andaloussi","Chiara Maggi","Alexandre Sollaci","Martin Stuermer","Petia Topalova"],"title":"Geoeconomic Fragmentation and Commodity Markets","year":2023,"url":"https://www.imf.org/-/media/files/publications/wp/2023/english/wpiea2023201-print-pdf.pdf","doi":"10.5089/9798400252426.001","locator":"Section 10, concluding research-agenda paragraph, PDF pages 39–40","evidence_summary":"Requests critical-mineral, food-security, recycling and elasticity extensions beyond simplifying scenarios."},"formalization_note":"The source explicitly requests refinement of its simplifying assumptions, food-security incidence, critical-mineral transition links, recycling and elasticities. The dynamic general-equilibrium scarcity-risk specification is proposed; its partial-equilibrium scenario results are already established conditional answers."},"statusQualification":"Published question or research agenda verified at the cited locator. The mathematical specification is adapted for this catalogue and includes additional assumptions; source publication does not establish first-ever priority or certify this exact model remains unsolved. The source explicitly requests refinement of its simplifying assumptions, food-security incidence, critical-mineral transition links, recycling and elasticities. The dynamic general-equilibrium scarcity-risk specification is proposed; its partial-equilibrium scenario results are already established conditional answers.","statusReviewedAt":"2026-10-08"}
% FIRST_POSTED: 2026-10-08
% ENTRY_KIND: statement_only
% !TeX program = lualatex
\documentclass[11pt]{article}
\usepackage{fontspec}
\usepackage[a4paper,margin=25mm]{geometry}
\usepackage{amsmath,amssymb,mathtools}
\usepackage{xurl}
\usepackage[hidelinks]{hyperref}
\newcommand{\sref}[2]{\hyperlink{source:#1:#2}{[S#2]}}
\setlength{\emergencystretch}{3em}
\begin{document}
% problemId:30007630
\section*{Commodity-trade fragmentation}\label{30007630}
\subsection{Definitions and mathematical statement}
Fix countries and commodity markets for food, energy and critical minerals, with production capacities, inventories and substitution technologies. A trade-fragmentation intervention \(a\) restricts specified bilateral flows at stated dates. Prices, investment and consumption adjust in a declared general-equilibrium system. Let \(p_{kct}(a)\) be the real price of commodity \(k\) in country \(c\), \(q_{kct}(a)\) consumption and \(W_c(a)\) household welfare. The task is to identify or bound dynamic changes in these outcomes and scarcity risk,
\[\Pr(q_{kc,t+h}(do(a))<\underline q_{kc}),\]
where \(\underline q_{kc}\) is an independently justified minimum requirement. Define the baseline policy path, commodity-grade aggregation and treatment of informal rerouting. Estimate substitution elasticities and supply-expansion lags rather than imposing instantaneous adaptation. Include links between energy prices, fertilizer production, food supply and mineral-intensive investment where empirically material. Distinguish temporary price redistribution from long-run productive-resource losses and evaluate impacts on countries with weak buffers. Report uncertainty from concentrated suppliers and correlated disruptions. Scenario results conditional on assumed blocs and elasticities do not establish the exact effects of an unspecified real-world fragmentation process.

\par\textbf{Formulation and scope.} The source explicitly requests refinement of its simplifying assumptions, food-security incidence, critical-mineral transition links, recycling and elasticities. The dynamic general-equilibrium scarcity-risk specification is proposed; its partial-equilibrium scenario results are already established conditional answers.

\par\textbf{Open-question provenance.} An explicit question or research agenda was verified at the source locator. This does not by itself certify that every later version remains unresolved \sref{30007630}{1}.

\par\textbf{Historical priority.} Earliest verified explicit extension agenda located in this audit; no first-statement priority claim for commodity fragmentation.

\subsection{Short English statement}
In the explicitly specified setting: How do restrictions on food, energy and critical-mineral trade change scarcity, substitution and development risks? The mathematical statement above fixes the target, assumptions and scope of this catalogue formulation.

\subsection{Sources}
\begin{itemize}\raggedright
\item[\textnormal{[S1]}]\hypertarget{source:30007630:1}{} Jorge A. Alvarez, Mehdi Benatiya Andaloussi, Chiara Maggi, Alexandre Sollaci, Martin Stuermer, Petia Topalova. 2023. Geoeconomic Fragmentation and Commodity Markets. Section 10, concluding research-agenda paragraph, PDF pages 39–40.
\url{https://www.imf.org/-/media/files/publications/wp/2023/english/wpiea2023201-print-pdf.pdf}.
DOI: 10.5089/9798400252426.001.
{\small\textit{Earliest verified question/agenda reference; historical priority qualified below. Requests critical-mineral, food-security, recycling and elasticity extensions beyond simplifying scenarios.}}
\end{itemize}
\end{document}
